\chapter{강화학습의 최소 토대: sleep을 하나의 의사결정 문제로 보기}
\label{bg:chapter-rl-foundations}

강화학습(Reinforcement Learning, RL)은 ``정답 label을 바로 주기 어렵지만 행동 뒤 결과는 관찰할 수
있는'' 순차 의사결정 문제를 다룬다. Sleep-time compute에서 RL의 주된 역할은 language model 전체를
무조건 RL로 다시 학습하는 것이 아니라, 언제 무엇을 저장·압축·승격할지 정하는 controller를 학습하는
데 있다.

\section{State, action, reward, trajectory}

\concept{rl-state}{RL 상태} $s_t$는 시점 $t$에서 정책이 볼 수 있는 정보다. 예를 들어 buffer fill ratio,
새 correction 수, retrieval miss, user inactivity window, available GPU, artifact age, privacy class가
state feature가 된다. state에 미래 정보나 evaluation answer가 들어가면 leakage다.

\concept{action}{행동} $a_t$는 ``아무것도 하지 않기'', summarize, re-embed, graph merge, replay batch
생성, LoRA training, candidate validate, promote, rollback처럼 시스템을 바꾸는 선택이다. action은
결과뿐 아니라 compute reservation, data read, model lock이라는 비용을 가진다.

\concept{reward}{보상} $r_t$는 행동의 즉시 결과다. 단순 accuracy gain만 쓰면 controller는 비싼 full
fine-tuning을 자주 고르거나, 쉬운 query를 암기해 점수를 올릴 수 있다. 그래서 reward를

\begin{equation}
r_t=\Delta U_t-\lambda_c C_t-\lambda_l L_t-\lambda_r Risk_t
-\lambda_f Forget_t
\end{equation}

처럼 utility 증가, compute/energy cost, wake latency, risk, forgetting으로 분해한다. 계수는 moral truth가
아니라 운영 정책이므로 버전과 승인 주체를 기록한다.

\concept{trajectory}{궤적} $\tau=(s_0,a_0,r_0,s_1,\ldots)$는 여러 시점의 순서다. 한 summary의 utility는
다음날 retrieval에서, 한 adapter의 부작용은 수주 뒤 rare slice에서 드러날 수 있다. 따라서 즉시 reward만
보면 장기 효과를 놓친다.

\section{Return, policy, value}

\concept{return}{리턴}은 미래 보상의 할인합이다.

\begin{equation}
G_t=\sum_{k=0}^{\infty}\gamma^k r_{t+k},\qquad 0\le\gamma<1.
\end{equation}

$\gamma$가 작으면 오늘의 compute 절약을, 크면 오래 뒤 retention과 maintenance를 더 본다. 삭제처럼
되돌리기 어려운 action에는 단일 discount factor보다 explicit safety constraint와 long-horizon audit가
필요하다.

\concept{policy}{정책} $\pi(a|s)$는 state에서 action distribution을 정한다. rule-based threshold도
policy이고 neural controller도 policy다. 먼저 rule baseline을 두어야 learned policy가 무엇을 개선했는지
알 수 있다. \concept{value-function}{가치 함수} $V^\pi(s)=\mathbb{E}[G_t|s_t=s]$는 state의 예상 장기
성과를, $Q^\pi(s,a)$는 action까지 정했을 때의 성과를 추정한다.

\concept{credit-assignment}{신용 할당}은 지연된 성공·실패를 어떤 과거 action에 돌릴지 정하는 문제다.
한 달 뒤 사용자가 오래된 preference를 정확히 회상했을 때 어느 nightly consolidation이 기여했는지
알기 어렵다. artifact lineage와 counterfactual shadow evaluation이 없으면 reward label 자체가 noisy해진다.

\section{Policy gradient와 update가 실제로 의미하는 것}

\concept{policy-gradient}{정책 기울기}는 기대 return
$J(\phi)=\mathbb{E}_{\tau\sim\pi_\phi}[G_0]$을 높이는 controller parameter $\phi$의 방향을 추정한다.

\begin{equation}
\nabla_\phi J(\phi)=\mathbb{E}\left[
\nabla_\phi\log\pi_\phi(a_t|s_t)\,A_t\right],
\end{equation}

$A_t$는 action이 평소 예상보다 얼마나 좋았는지 나타내는 advantage다. PPO는 한 batch를 여러 epoch
사용하되 policy가 한 번에 너무 멀리 움직이지 않도록 clipped surrogate objective를 사용한다
\parencite{schulman2017ppo}. 이름을 외우기보다 ``rollout 수집 → advantage 추정 → policy update →
새 policy로 다시 수집''이라는 dependency를 이해해야 한다.

\section{On-policy와 off-policy의 systems 차이}

\concept{on-policy}{온폴리시} 학습은 현재 policy가 만든 experience를 사용한다. distribution mismatch는
작지만 새 policy마다 실제 또는 simulator rollout이 필요하고, unsafe action을 exploration할 위험이 있다.
\concept{off-policy}{오프폴리시} 학습은 과거 policy나 rule controller가 만든 log도 사용한다. 기존 운영
trace를 재활용할 수 있지만, 과거에 거의 선택되지 않은 action의 결과는 데이터에 없다는 support 문제가
생긴다.

Sleep controller에서는 실제 user data로 무작정 exploration하지 않는다. 안전한 순서는 (1) trace-driven
simulator, (2) counterfactual cost model, (3) shadow execution, (4) 제한된 canary, (5) gradual rollout이다.
``do nothing''과 현재 heuristic은 항상 baseline action으로 남긴다.

\begin{table}[htbp]
\centering
\caption{지도학습과 강화학습을 sleep control에 적용할 때의 차이}
\label{tab:bg-supervised-vs-rl}
\small
\begin{tabularx}{\textwidth}{p{0.18\textwidth} X X}
\toprule
항목 & 지도/선호 분류 & 순차 RL \\
\midrule
label & 이 state에서 어떤 action이 좋은가 & action 뒤 여러 시점의 reward \\
data & 독립 예시로 근사 가능 & policy가 방문한 trajectory distribution \\
강점 & 안정적이고 offline 검증이 쉬움 & 지연 효과와 action sequence 최적화 가능 \\
위험 & expert label 비용, myopic target & exploration, reward hacking, counterfactual 부재 \\
권장 시작 & heuristic imitation, pairwise ranking & simulator와 shadow에서 제한적으로 \\
\bottomrule
\end{tabularx}
\end{table}

\begin{cautionbox}[Reward가 곧 요구사항은 아니다]
측정 가능한 proxy를 reward로 쓰면 controller는 proxy를 최적화한다. retrieval hit rate만 높이면 불필요한
기억을 모두 보존할 수 있고, compression ratio만 높이면 provenance를 지울 수 있다. hard constraint,
slice별 validation, human override, action audit를 reward 바깥에 둔다.
\end{cautionbox}
